\documentclass[12pt,a4paper]{article}
\usepackage[T2A]{fontenc}
\usepackage[utf8]{inputenc}
\usepackage[russian]{babel}
\usepackage[left=30mm,right=20mm,top=20mm,bottom=20mm]{geometry}
\usepackage{graphicx}
\usepackage{float}
\usepackage{array,longtable,tabularx}
\usepackage{listings}
\usepackage{xcolor}
\usepackage[hidelinks]{hyperref}
\setlength{\parindent}{1.25cm}
\setlength{\parskip}{4pt}
\renewcommand{\arraystretch}{1.2}
\emergencystretch=2em
\lstset{
  basicstyle=\ttfamily\small,
  breaklines=true,
  columns=fullflexible,
  keepspaces=true,
  showstringspaces=false,
  frame=single,
  rulecolor=\color{gray},
  tabsize=4
}
% Укажите имена файлов готовых диаграмм вместо conceptual.png и er.png.
% Если изображений пока нет рядом с report.tex, документ всё равно компилируется.
\newcommand{\readydiagram}[2]{%
  \begin{figure}[H]
    \centering
    \IfFileExists{#1}{%
      \includegraphics[width=\linewidth,height=0.65\textheight,keepaspectratio]{#1}%
    }{%
      \fbox{\parbox{0.9\linewidth}{\centering
        #2\\[6pt]
        Файл готовой диаграммы: \texttt{\detokenize{#1}}}}
    }
    \caption{#2}
  \end{figure}
}
\title{Отчёт по лабораторной работе 1\\Информационная система квартир}
\author{}
\date{2026}

\begin{document}
\maketitle

В работе рассматриваются проектирование и реализация базы данных информационной
системы для управления квартирами. Описаны предметная область, сущности,
концептуальная, логическая и физическая модели, а также SQL-схема и специальные
операции. Реализация использует PostgreSQL, Jakarta EE и JPA с EclipseLink.

\section{Текст задания}

Реализовать информационную систему, которая позволяет взаимодействовать с
объектами класса \texttt{Flat}. Основной объект содержит следующие поля:
\texttt{id}, \texttt{name}, \texttt{coordinates}, \texttt{creationDate},
\texttt{area}, \texttt{price}, \texttt{balcony}, \texttt{timeToMetroOnFoot},
\texttt{numberOfRooms}, \texttt{furniture}, \texttt{view}, \texttt{transport},
\texttt{house}.

Требования к полям объекта и вспомогательных классов:
\begin{itemize}
  \item \texttt{Flat.id}: положительный уникальный идентификатор,
  генерируемый автоматически средствами БД.
  \item \texttt{Flat.name}: непустая строка, не равная \texttt{null}.
  \item \texttt{Flat.coordinates}: обязательный объект координат.
  \item \texttt{Flat.creationDate}: обязательная дата и время создания,
  формируемые автоматически.
  \item \texttt{Flat.area}, \texttt{Flat.price},
  \texttt{Flat.timeToMetroOnFoot}: положительные значения.
  \item \texttt{Flat.numberOfRooms}: целое число от 1 до 11.
  \item \texttt{Flat.balcony}, \texttt{Flat.furniture}: значения
  \texttt{Boolean}, допускающие \texttt{null}.
  \item \texttt{Flat.view}: обязательное значение перечисления
  \texttt{View}: \texttt{STREET}, \texttt{YARD}, \texttt{PARK},
  \texttt{BAD}, \texttt{TERRIBLE}.
  \item \texttt{Flat.transport}: необязательное значение перечисления
  \texttt{Transport}: \texttt{FEW}, \texttt{LITTLE}, \texttt{NORMAL}.
  \item \texttt{Flat.house}: необязательный объект дома.
  \item \texttt{Coordinates.x}, \texttt{Coordinates.y}: обязательные значения
  типа \texttt{Float}; значение \texttt{y} не превышает 249.
  \item \texttt{House.name}: строка, не равная \texttt{null}.
  \item \texttt{House.year}: обязательное положительное значение типа
  \texttt{Long}; \texttt{House.numberOfFloors}: положительное целое число.
\end{itemize}

Разработанная система должна удовлетворять следующим требованиям:
\begin{enumerate}
  \item Основное назначение системы --- управление объектами заданного класса.
  \item Предусмотреть создание, получение информации по ID, изменение и
  удаление объектов. Операции выполняются в отдельных окнах или интерфейсах.
  При просмотре объекта отображается информация о связанных объектах.
  \item При создании квартиры обеспечить выбор существующих объектов
  вспомогательных классов.
  \item Выполнять операции управления на сервере и синхронизировать изменения с БД.
  \item На главном экране отображать таблицу объектов, выделяя каждому атрибуту
  отдельную колонку. При необходимости использовать пагинацию.
  \item Обеспечить сортировку и фильтрацию по значениям строковых колонок.
  Фильтрация должна выполняться по неполному совпадению.
  \item Предоставить переход к изменению объекта из общей таблицы и из области
  визуализации объекта, если она реализована.
  \item При добавлении, изменении и удалении объектов автоматически обновлять
  интерфейсы других авторизованных пользователей.
  \item Запрещать удаление объекта, на который ссылается другой объект,
  и сообщать пользователю о невозможности удаления.
  \item Разрешить пользователям просмотр всех объектов. Для изменения
  открывать отдельное диалоговое окно. При некорректном вводе показывать
  информативные сообщения.
  \item Предусмотреть меню для перехода между логически обособленными частями системы.
\end{enumerate}

В отдельном интерфейсе должны быть доступны пять специальных операций:
\begin{enumerate}
  \item Вернуть один объект с максимальным значением \texttt{area}.
  \item Вернуть количество квартир, значение \texttt{house} которых равно заданному.
  \item Вернуть массив квартир, название которых начинается с заданной подстроки.
  \item Найти самую дорогую квартиру без балкона.
  \item Выбрать наиболее дорогую из трёх квартир с заданными ID.
\end{enumerate}
Операции необходимо реализовать как функции БД и вызывать из бизнес-логики.

Хранение данных организовать в PostgreSQL. Ограничения полей должны быть
реализованы на уровне ORM и БД. Для кафедрального развёртывания использовать
хост \texttt{pg}, базу \texttt{studs} и учётные данные пользователя кафедрального
сервера. Основой приложения должны служить Java EE и управляемые компоненты,
уровнем хранения --- JPA с EclipseLink. Уровни и логические части приложения
необходимо разделить на отдельные компоненты.

\section{Описание предметной области}

Предметная область --- учёт квартир и их характеристик. Для квартиры хранятся
название, площадь, цена, число комнат, время пешком до метро, сведения о балконе
и мебели, вид из окна и доступность транспорта. Местоположение задаётся
координатами, а сведения о здании выделены в объект дома.

Квартира обязательно связана с координатами. Связь с домом необязательна:
допускается запись квартиры без сведений о доме. Один дом и один объект
координат могут использоваться несколькими квартирами. Дом содержит название,
год и количество этажей; координаты --- значения \texttt{x} и \texttt{y}.

Зарегистрированный пользователь просматривает квартиры, выбирает связанные
объекты и выполняет операции управления. Учётная запись содержит логин и хеш
пароля. В модели не предусмотрена принадлежность квартиры пользователю:
все авторизованные пользователи работают с общим набором данных.

При изменении квартиры сервер сохраняет данные в транзакции. После успешного
коммита через WebSocket отправляется уведомление, по которому интерфейсы
повторно запрашивают актуальные данные через REST. Прямое изменение БД
внешним SQL-запросом само по себе такое уведомление не создаёт.

\section{Сущности и концептуальная модель}

\subsection{Список и классификация сущностей}

\begin{tabularx}{\linewidth}{|p{2.5cm}|p{3.8cm}|X|}
\hline
\textbf{Сущность} & \textbf{Классификация} & \textbf{Назначение} \\ \hline
Квартира & Основная предметная, сильная & Объект учёта, содержащий характеристики жилья и ссылки на связанные объекты. \\ \hline
Дом & Вспомогательная предметная, сильная & Общие сведения о здании, которые могут использоваться несколькими квартирами. \\ \hline
Координаты & Вспомогательная предметная, сильная & Местоположение; самостоятельная запись с собственным идентификатором. \\ \hline
Пользователь & Служебная, сильная & Учётная запись для регистрации и входа в систему. \\ \hline
\end{tabularx}

Все четыре сущности имеют собственные первичные ключи, поэтому в данной модели
нет слабых сущностей и таблиц для связей многие ко многим. Обязательная связь
квартиры с координатами не делает квартиру слабой сущностью: её идентификатор
не зависит от идентификатора координат.

\texttt{View} и \texttt{Transport} --- перечислимые домены значений атрибутов,
а не самостоятельные сущности. DTO, сервисы, контроллеры, JWT и WebSocket-сессии
также не являются сущностями базы данных.

\subsection{Связи концептуальной модели}

\begin{itemize}
  \item Дом --- квартира: один дом связан с нулём или несколькими квартирами;
  квартира связана с нулём или одним домом.
  \item Координаты --- квартира: один объект координат связан с нулём или
  несколькими квартирами; квартира связана ровно с одним объектом координат.
  \item Пользователь не имеет хранимой связи с квартирой, домом или координатами.
  Его участие в операциях системы не моделируется внешним ключом.
\end{itemize}

\readydiagram{conceptual.png}{Концептуальная модель предметной области}

\section{Логическая модель и ER-диаграмма}

Концептуальные сущности преобразованы в четыре отношения. В каждой таблице
используется суррогатный первичный ключ \texttt{id}.

\begin{itemize}
  \item \texttt{coordinates}: \texttt{id} (PK), \texttt{x}, \texttt{y}.
  \item \texttt{house}: \texttt{id} (PK), \texttt{name}, \texttt{year},
  \texttt{number\_of\_floors}.
  \item \texttt{flat}: \texttt{id} (PK), \texttt{name},
  \texttt{coordinates\_id} (FK), \texttt{creation\_date}, \texttt{area},
  \texttt{price}, \texttt{balcony}, \texttt{time\_to\_metro\_on\_foot},
  \texttt{number\_of\_rooms}, \texttt{furniture}, \texttt{view\_type},
  \texttt{transport\_type}, \texttt{house\_id} (FK).
  \item \texttt{app\_user}: \texttt{id} (PK), \texttt{username} (уникальный ключ),
  \texttt{password\_hash}.
\end{itemize}

\texttt{flat.coordinates\_id} ссылается на \texttt{coordinates.id} и не
допускает \texttt{NULL}. \texttt{flat.house\_id} ссылается на \texttt{house.id}
и допускает \texttt{NULL}. Обе связи имеют тип один ко многим.

\readydiagram{er.png}{Логическая модель базы данных и ER-диаграмма}

Атрибуты содержат атомарные значения, повторяющиеся группы отсутствуют.
Первичные ключи состоят из одного атрибута, поэтому частичных зависимостей
от составного ключа нет. Характеристики дома и координаты вынесены из
таблицы квартир в отдельные отношения. При принятых зависимостях от ключей
модель соответствует третьей нормальной форме.

Полный ответ API может содержать название дома, год, этажность и координаты
в одном DTO. Это объединение данных для отображения, а не дублирование
этих полей в таблице \texttt{flat}.

\section{Физическая модель}

СУБД --- PostgreSQL 17. Прикладные идентификаторы генерируются с помощью
\texttt{GENERATED ALWAYS AS IDENTITY}. В JPA используется стратегия
\texttt{GenerationType.IDENTITY}. Java-классы отображаются на таблицы
аннотациями \texttt{@Entity}, \texttt{@Table}, \texttt{@Column} и
\texttt{@JoinColumn}. Автоматическое создание схемы ORM отключено:
структуру БД задают SQL-скрипты.

В таблицах ниже обязательность означает наличие ограничения \texttt{NOT NULL}.

\subsection{Таблица coordinates}
\begin{tabularx}{\linewidth}{|p{1.6cm}|p{2.1cm}|p{2.6cm}|X|}
\hline
\textbf{Поле} & \textbf{Тип} & \textbf{Обязательность} & \textbf{Ограничения} \\ \hline
\texttt{id} & BIGINT & Обязательно & Первичный ключ, IDENTITY, больше нуля. \\ \hline
\texttt{x} & REAL & Обязательно & Ограничение диапазона в SQL не задано. \\ \hline
\texttt{y} & REAL & Обязательно & Значение не превышает 249. \\ \hline
\end{tabularx}

\subsection{Таблица house}
\begin{tabularx}{\linewidth}{|p{3.4cm}|p{2cm}|p{2.6cm}|X|}
\hline
\textbf{Поле} & \textbf{Тип} & \textbf{Обязательность} & \textbf{Ограничения} \\ \hline
\texttt{id} & BIGINT & Обязательно & Первичный ключ, IDENTITY, больше нуля. \\ \hline
\texttt{name} & TEXT & Обязательно & Пустая строка допускается. \\ \hline
\texttt{year} & BIGINT & Обязательно & Больше нуля. \\ \hline
\texttt{number\_of\_floors} & INTEGER & Обязательно & Больше нуля. \\ \hline
\end{tabularx}

\subsection{Таблица flat}
\begin{longtable}{|p{4.8cm}|p{2.5cm}|p{1.6cm}|p{4.3cm}|}
\hline
\textbf{Поле} & \textbf{Тип} & \textbf{NULL} & \textbf{Ограничения} \\ \hline
\endfirsthead
\hline
\textbf{Поле} & \textbf{Тип} & \textbf{NULL} & \textbf{Ограничения} \\ \hline
\endhead
\texttt{id} & INTEGER & Нет & Первичный ключ, IDENTITY, больше нуля. \\ \hline
\texttt{name} & TEXT & Нет & Содержит хотя бы один непробельный символ. \\ \hline
\texttt{coordinates\_id} & BIGINT & Нет & Внешний ключ на \texttt{coordinates.id}. \\ \hline
\texttt{creation\_date} & TIMESTAMP WITHOUT TIME ZONE & Нет & По умолчанию \texttt{CURRENT\_TIMESTAMP}. \\ \hline
\texttt{area} & BIGINT & Нет & Больше нуля. \\ \hline
\texttt{price} & DOUBLE PRECISION & Нет & Конечное значение больше нуля. \\ \hline
\texttt{balcony} & BOOLEAN & Да & \texttt{TRUE}, \texttt{FALSE} или \texttt{NULL}. \\ \hline
\texttt{time\_to\_metro\_on\_foot} & REAL & Нет & Конечное значение больше нуля. \\ \hline
\texttt{number\_of\_rooms} & INTEGER & Нет & От 1 до 11 включительно. \\ \hline
\texttt{furniture} & BOOLEAN & Да & \texttt{TRUE}, \texttt{FALSE} или \texttt{NULL}. \\ \hline
\texttt{view\_type} & VARCHAR(16) & Нет & Значение из перечисления \texttt{View}. \\ \hline
\texttt{transport\_type} & VARCHAR(16) & Да & Значение из \texttt{Transport} или \texttt{NULL}. \\ \hline
\texttt{house\_id} & BIGINT & Да & Внешний ключ на \texttt{house.id}. \\ \hline
\end{longtable}

Дата создания формируется в JPA-методе с \texttt{@PrePersist}. SQL-значение
по умолчанию используется при вставке без явного указания даты. ORM исключает
дату из обновлений через \texttt{updatable=false}; отдельного SQL-триггера,
запрещающего изменение даты прямым запросом, нет.

\subsection{Таблица app\_user}
\begin{tabularx}{\linewidth}{|p{3.2cm}|p{2.4cm}|p{2.6cm}|X|}
\hline
\textbf{Поле} & \textbf{Тип} & \textbf{Обязательность} & \textbf{Ограничения} \\ \hline
\texttt{id} & BIGINT & Обязательно & Первичный ключ, IDENTITY. \\ \hline
\texttt{username} & VARCHAR(32) & Обязательно & Уникальный логин: 3--32 строчных латинских буквы, цифры или символы подчёркивания. \\ \hline
\texttt{password\_hash} & VARCHAR(200) & Обязательно & Хеш пароля; открытый пароль не хранится. \\ \hline
\end{tabularx}

Пароли хешируются алгоритмом PBKDF2-HMAC-SHA256 с индивидуальной случайной
солью и 600\,000 итерациями. Параметры алгоритма, соль и результат хеширования
сохраняются одной строкой. JWT выдаётся после регистрации или входа и не
хранится отдельной записью в БД.

\subsection{Целостность и индексы}

Первичные ключи обеспечивают идентификацию строк. \texttt{NOT NULL} задаёт
обязательность полей, \texttt{CHECK} ограничивает диапазоны и домены, а
\texttt{UNIQUE} запрещает повторение логинов. Перечисления сохраняются
строками, для которых предусмотрены SQL-проверки допустимых значений.

Внешние ключи используют \texttt{ON DELETE RESTRICT}: нельзя удалить дом
или координаты, пока на них ссылаются квартиры. Удаление квартиры не удаляет
связанные записи. Каскадное удаление не настроено; API удаления домов и
координат в текущей реализации отсутствует.

PostgreSQL создаёт индексы для первичных ключей и уникального логина.
Дополнительные индексы для внешних ключей, сортировки и поиска по подстроке
в SQL-схеме не объявлены.

\section{Реализация модели на SQL}

Схема определена в \texttt{db/init/schema.sql}, учётные записи --- в
\texttt{db/init/users.sql}, специальные операции --- в
\texttt{db/init/flat\_functions.sql}. Ниже приведены определения из этих файлов.

\subsection{Создание предметных таблиц}
\begin{lstlisting}[language=SQL]
BEGIN;

CREATE TABLE coordinates (
    id BIGINT GENERATED ALWAYS AS IDENTITY PRIMARY KEY,
    x REAL NOT NULL,
    y REAL NOT NULL,
    CONSTRAINT ck_coordinates_id_positive CHECK (id > 0),
    CONSTRAINT ck_coordinates_y_max CHECK (y <= 249)
);

CREATE TABLE house (
    id BIGINT GENERATED ALWAYS AS IDENTITY PRIMARY KEY,
    name TEXT NOT NULL,
    year BIGINT NOT NULL,
    number_of_floors INTEGER NOT NULL,
    CONSTRAINT ck_house_id_positive CHECK (id > 0),
    CONSTRAINT ck_house_year_positive CHECK (year > 0),
    CONSTRAINT ck_house_floors_positive CHECK (number_of_floors > 0)
);

CREATE TABLE flat (
    id INTEGER GENERATED ALWAYS AS IDENTITY PRIMARY KEY,
    name TEXT NOT NULL,
    coordinates_id BIGINT NOT NULL,
    creation_date TIMESTAMP WITHOUT TIME ZONE NOT NULL
        DEFAULT CURRENT_TIMESTAMP,
    area BIGINT NOT NULL,
    price DOUBLE PRECISION NOT NULL,
    balcony BOOLEAN,
    time_to_metro_on_foot REAL NOT NULL,
    number_of_rooms INTEGER NOT NULL,
    furniture BOOLEAN,
    view_type VARCHAR(16) NOT NULL,
    transport_type VARCHAR(16),
    house_id BIGINT,
    CONSTRAINT ck_flat_id_positive CHECK (id > 0),
    CONSTRAINT ck_flat_name_not_blank CHECK (name ~ '[^[:space:]]'),
    CONSTRAINT ck_flat_area_positive CHECK (area > 0),
    CONSTRAINT ck_flat_price_positive_finite
        CHECK (price > 0 AND price < 'Infinity'::DOUBLE PRECISION),
    CONSTRAINT ck_flat_metro_time_positive_finite
        CHECK (time_to_metro_on_foot > 0
            AND time_to_metro_on_foot < 'Infinity'::REAL),
    CONSTRAINT ck_flat_rooms_range
        CHECK (number_of_rooms BETWEEN 1 AND 11),
    CONSTRAINT ck_flat_view
        CHECK (view_type IN ('STREET', 'YARD', 'PARK', 'BAD', 'TERRIBLE')),
    CONSTRAINT ck_flat_transport
        CHECK (transport_type IS NULL
            OR transport_type IN ('FEW', 'LITTLE', 'NORMAL')),
    CONSTRAINT fk_flat_coordinates
        FOREIGN KEY (coordinates_id) REFERENCES coordinates(id)
        ON DELETE RESTRICT,
    CONSTRAINT fk_flat_house
        FOREIGN KEY (house_id) REFERENCES house(id)
        ON DELETE RESTRICT
);

COMMIT;
\end{lstlisting}

\subsection{Создание таблицы пользователей}
\begin{lstlisting}[language=SQL]
BEGIN;
CREATE TABLE IF NOT EXISTS app_user (
    id BIGINT GENERATED ALWAYS AS IDENTITY PRIMARY KEY,
    username VARCHAR(32) NOT NULL UNIQUE,
    password_hash VARCHAR(200) NOT NULL,
    CONSTRAINT ck_user_username
        CHECK (username ~ '^[a-z0-9_]{3,32}$')
);
COMMIT;
\end{lstlisting}

\subsection{Специальные операции}

Четыре функции реализованы на SQL, выбор самой дорогой из трёх квартир ---
на PL/pgSQL. При равных площади или цене выбирается объект с наименьшим ID.
\texttt{STABLE} характеризует поведение функций в пределах SQL-запроса
и не означает постоянного кеширования результатов.

\begin{lstlisting}[language=SQL]
BEGIN;

CREATE OR REPLACE FUNCTION flat_max_area()
RETURNS SETOF flat
LANGUAGE sql STABLE
AS $$
    SELECT f.* FROM flat f
    ORDER BY f.area DESC, f.id ASC LIMIT 1;
$$;

CREATE OR REPLACE FUNCTION flat_count_by_house(p_house_id BIGINT)
RETURNS BIGINT
LANGUAGE sql STABLE
AS $$
    SELECT count(*) FROM flat f
    WHERE f.house_id IS NOT DISTINCT FROM p_house_id;
$$;

CREATE OR REPLACE FUNCTION flat_names_starting_with(p_prefix TEXT)
RETURNS SETOF flat
LANGUAGE sql STABLE
AS $$
    SELECT f.* FROM flat f
    WHERE left(f.name, char_length(p_prefix)) = p_prefix
    ORDER BY f.id ASC;
$$;

CREATE OR REPLACE FUNCTION flat_most_expensive_without_balcony()
RETURNS SETOF flat
LANGUAGE sql STABLE
AS $$
    SELECT f.* FROM flat f WHERE f.balcony IS FALSE
    ORDER BY f.price DESC, f.id ASC LIMIT 1;
$$;

CREATE OR REPLACE FUNCTION flat_most_expensive_of_three(
    p_id1 INTEGER, p_id2 INTEGER, p_id3 INTEGER)
RETURNS SETOF flat
LANGUAGE plpgsql STABLE
AS $$
BEGIN
    RETURN QUERY
        SELECT f.* FROM flat f
        WHERE f.id IN (p_id1, p_id2, p_id3)
        ORDER BY f.price DESC, f.id ASC LIMIT 1;
END;
$$;

COMMIT;
\end{lstlisting}

\texttt{flat\_count\_by\_house(NULL)} возвращает количество квартир без дома.
Для отсутствующего ID дома результат равен нулю. Поиск по префиксу учитывает
регистр; пустой префикс соответствует всем названиям. В операции поиска
без балкона учитываются только строки с \texttt{balcony=FALSE}; неизвестное
значение \texttt{NULL} не считается отсутствием балкона.

Проверка трёх положительных, различных и существующих ID выполняется в
\texttt{FlatOperationsService}. Сама SQL-функция выбирает максимум среди
найденных строк и не содержит этих проверок.

Примеры вызовов функций:
\begin{lstlisting}[language=SQL]
SELECT * FROM flat_max_area();
SELECT flat_count_by_house(1);
SELECT flat_count_by_house(NULL);
SELECT * FROM flat_names_starting_with('Studio');
SELECT * FROM flat_most_expensive_without_balcony();
SELECT * FROM flat_most_expensive_of_three(1, 2, 3);
\end{lstlisting}

\subsection{Связь SQL-модели с приложением}

REST-контроллеры вызывают EJB-сервисы; сервисы выполняют проверки и обращаются
к репозиториям. Репозитории используют \texttt{EntityManager}. Списки квартир,
домов и координат строятся через Criteria API; специальные операции вызываются
через параметризованные native queries. При регистрации применяется
\texttt{INSERT ... ON CONFLICT DO NOTHING}, чтобы уникальное ограничение
обрабатывало в том числе одновременные регистрации одного логина.

EJB управляет транзакциями. \texttt{flush()} синхронизирует изменения с БД,
но не заменяет коммит. Уведомление об изменении отправляется только после
успешного завершения транзакции.

При первом запуске PostgreSQL в Docker таблицы и функции создаются init-скриптами.
Для существующей БД новые определения применяются отдельно; пересборка WAR
не изменяет сохранённую схему автоматически.

\section{Выводы по работе}

В ходе работы спроектирована и реализована реляционная модель информационной
системы квартир. Выделены четыре сущности, определены первичные и внешние ключи,
обязательность атрибутов и ограничения допустимых значений. Разделение квартир,
домов и координат позволяет повторно использовать связанные записи и хранить
общие характеристики без дублирования.

Полученная модель реализована в PostgreSQL и связана с Java-приложением через
JPA и EclipseLink. Операции изменения выполняются в транзакциях, пять
специальных операций реализованы функциями БД. Регистрация и вход используют
хеширование паролей и JWT; уведомления WebSocket позволяют обновлять интерфейсы
после сохранения изменений.

Работа позволила применить на практике переход от предметной области к
концептуальной, логической и физической моделям, ограничения ссылочной
целостности и интеграцию ORM с SQL-функциями.

\end{document}
